\documentclass[11pt]{article}
\usepackage{amsmath,amssymb}
\usepackage[margin=1in]{geometry}
\usepackage{hyperref}

\title{Ideas Parked --- Not Yet Merged (v18)}
\author{Staging document; each entry marks its eventual destination}
\date{}

\begin{document}
\maketitle

\noindent The full working-out of Gap~3 (ideas\_v1: ``reducing a
path by merging close actors''), reached over several turns and
recorded here as one connected argument. Its ending differs from
Gap~1's and Gap~2's: this is not a resolution to merge into Paper~1,
but a determination that most of the question was never Paper~1's
to answer. Paper~1 has already been narrowed accordingly
(Pattern\_Arithmetic\_v15\_22), keeping only the mechanical
guarantee; everything below is destined for Inference instead.

\section{Where the discussion started, and the correction made along
the way}

The first pass treated ``direction'' in Gap~3's own notation
($\chi'=\arg(z)$) as $\hat n$, and invoked Gap~1's refusal of Mix
across differing $\hat n$ against it. That was wrong, and worth
recording as a correction rather than quietly dropped: $\chi$ is
phase at a \emph{shared} direction, the same coordinate Mix's own
flagship example combines --- red ($\chi=0°$) plus green
($\chi=120°$) gives yellow ($\chi=60°$), verified directly,
Mix's validated central case, not a refused one. So combining
different $\chi$ values at one shared $\hat n$ was never the
problem. The real issue, once that confusion was cleared, turned out
to be order-blindness: checked directly, three actors merged in one
sequence and the same three merged in reverse give identical
magnitude and angle, since addition is commutative regardless of
order.

\section{Where a path's commitment to order actually comes from}

Checked against the paper's own words: an observed sequence arrives
already ordered from outside the system; Bundle's first act is to
discard that order deliberately ($A\uplus A=A$, order-blind by
construction); $T$'s actors are then ``taken from the bowl or from
the tape'' --- a choice made when a walk is built, not generated by
any axiom. The sharp consequence: choosing to route data through $T$
rather than leaving it in a bowl \emph{is itself} the declaration
that order matters for that data. Nothing mints that commitment
internally; the modeler asserts it once, by which operator they
reach for.

\section{The temperature example, and why it was never a
counterexample}

10,000 readings in a minute, reduced to one average: checked against
the above, this is Bowl $\to$ Mix, never routed through $T$ at all.
The decision ``one reading per minute suffices'' \emph{is} the
declaration that intra-minute order is noise. Order-blindness only
bites data that was already committed to $T$ --- asserting, by that
earlier choice, that the trajectory itself was part of the answer.
The same raw readings support at least three different questions,
and which representation is correct depends on which question is
being asked, not on the data itself:
\begin{itemize}
\item \emph{What was the average?} --- order within the window is
noise. Bowl, Mix, one point. Correct as given.
\item \emph{Did it spike and return?} --- the trajectory is the
answer; collapsing to one point erases the one thing being asked
for.
\item \emph{Is the broader trend rising or falling?} --- now the
sequence of per-window averages is itself ordered, and that order
matters again, one level up.
\end{itemize}

\section{Three techniques, verified as answering three different
questions, not one technique with three names}

\paragraph{Average (one point, possibly manufactured).} Checked
directly: a spike path ($10°\!\to\!40°\!\to\!70°\!\to\!40°\!\to\!10°$)
and a flat path sitting at $40°$ throughout produce the identical
single output ($33.79°$ averaged, or $40°$ for the real closest
point either way) --- one number cannot distinguish ``it spiked''
from ``it never moved,'' whether that number is manufactured or
picked from the real data. Realness was never the issue; count was.

\paragraph{Multiple centroids (several points, grouped by value,
blind to position).} Checked directly: a forward spike sequence and
its reverse --- same multiset of values, opposite order --- produce
\emph{identical} $2$-means centroids ($\{22,70\}$ either way).
Clustering answers ``what regimes occurred'' correctly and
efficiently; it cannot answer ``in what order,'' by construction,
since position was never an input to the grouping.

\paragraph{Shape-preserving reduction (several real points, kept at
their original positions).} Keeping start/peak/end from the same
spike --- three real values, in original sequence --- lets a reader
see the rise and the return. This is not a new technique: it is
exactly what line-simplification algorithms (Douglas--Peucker,
canonically) already do for ordinary paths, parameterized by a
tolerance $\varepsilon$ that bounds how far the simplified path may
deviate from the original. That $\varepsilon$ \emph{is} the drift
bound ideas\_v1's own Gap~3 text asked for, under a different name,
already solved elsewhere, not yet ported here.

\section{What stays in Paper 1, and why (already executed)}

Per the algebra's own existing principle --- a reduction is a
second, optional reading, not a replacement, the same promise
Collapse already keeps for its own sources --- Paper~1 states only
that any reduction must not consume what it reduces. It does not
choose among the three techniques above, because that choice depends
on purpose, and purpose is Inference's business, not the algebra's.
This has already been written into
\texttt{Pattern\_Arithmetic\_v15\_22.tex}, narrowing the former
Future-Work paragraph to exactly this mechanical content and
redirecting the rest here.

\section{What is destined for Inference, specifically}

\begin{itemize}
\item The technique-selection question itself: average versus
centroid versus shape-preservation, chosen by what the reduction
needs to answer, as argued throughout this entry.
\item The stated distance and drift bound, meaningful only once a
technique and a purpose are both fixed --- not general algebra
facts, but parameters of a specific reference-atom or classification
need.
\item The one genuinely unbuilt piece, narrowed considerably by this
discussion: adapting Douglas--Peucker's planar deviation measure to
this geometry specifically. The original algorithm measures
straight-line distance from a chord between two kept points; here,
``the chord between two kept actors'' needs to mean something
stated in terms of $T$'s own rotations --- plausibly geodesic
deviation from the great-circle arc the two kept actors would
define, not planar distance. This adaptation, not the general
principle, is the remaining work.
\end{itemize}

\section{Status}

Gap~3 is closed in the sense that matters for Paper~1: the paper no
longer owes an answer it was never positioned to give, and says so
directly rather than leaving an unresolved paragraph sitting in its
own Future Work. The technique-selection question is fully specified
here --- three distinct tools, each verified against the question it
actually answers, not three names for one idea --- and ready for
Inference to pick up, pending the one adaptation named above.

\end{document}
